\documentclass{midl} 

\usepackage{mwe} % to get dummy images
\jmlrvolume{}
\jmlryear{2026}
\jmlrworkshop{MVA ENS Paris-Saclay - CentraleSupélec} 
\editors{}

\title[Histopathology OOD Classification]{MVA DLMI 2026 - Histopathology OOD Classification} 

\midlauthor{\Name{Ayoub El Kbadi\midljointauthortext{Contributed equally}\nametag{$^{1,2}$}} \Email{ayoub.kbdi@gmail.com}\\
\addr $^{1}$ CentraleSupélec \\
\addr $^{2}$ ENS Paris-Saclay \AND  
\Name{Jean-Vincent Martini\midlotherjointauthor\nametag{$^{1,2}$}} \Email{jean-vincent.martini@ens-paris-saclay.fr}\\
\addr $^{1}$ CentraleSupélec \\
\addr $^{2}$ ENS Paris-Saclay
} 

\begin{document}

\maketitle

\begin{abstract}
This is a great paper and it has a concise abstract.
\end{abstract}

\begin{keywords}
List of keywords, comma separated.
\end{keywords}

\section{Introduction} 
TODO
 
\section{Problem Formulation \& Dataset}
 
\section{Proposed Approach}
  \subsection{Baseline}
\subsection{TODO: our methods}

\section{Experiments \& Results}

\section{Conclusion}

\begin{table}[htbp]
\floatconts
  {tab:example}%
  {\caption{An Example Table}}%
  {\begin{tabular}{ll}
  \bfseries Dataset & \bfseries Result\\
  Data1 & 0.12345\\
  Data2 & 0.67890\\
  Data3 & 0.54321\\
  Data4 & 0.09876
  \end{tabular}}
\end{table}

\begin{figure}[htbp]
\floatconts
  {fig:example}
  {\caption{Example Image}}
  {\includegraphics[width=0.5\linewidth]{example-image}}
\end{figure}

\begin{algorithm2e}
\caption{Random example}
\label{alg:net}
\KwIn{$x_1, \ldots, x_n, w_1, \ldots, w_n$}
\KwOut{$y$, the net activation}
$y\leftarrow 0$\;
\For{$i\leftarrow 1$ \KwTo $n$}{
  $y \leftarrow y + w_i*x_i$\;
}
\end{algorithm2e}

\clearpage  

\midlacknowledgments{We thank a bunch of people.}


\bibliography{bibliography}


\appendix

\section{Appendix}
TODO

\end{document}
